\iffalse
	\title{2023}
	\author{AI24BTECH11003}
	\section{mcq-single}
\fi
 
%16
    \item If the tangent at a point $P$ on the parabola $y^2 = 3x$ is parallel to the line $x+2y=1$ and the tangents at the points $Q$ and $R$ on the ellipse $\frac{x^2}{4}+\frac{y^2}{1}=1$ are perpendicular to the line $x-y=2$, then the area of the triangle $PQR$ is:
    
    \hfill{[Jan 2023]}

        \begin{multicols}{4}
            \begin{enumerate}
                \item $\frac{9}{\sqrt{5}}$
                \item $5\sqrt{3}$
                \item $\frac{3}{2}\sqrt{5}$
                \item $3\sqrt{5}$
            \end{enumerate}
        \end{multicols}

%17
    \item Let $y=y\brak{x}$ be the solution of the differential equation $x\log_ex\frac{dy}{dx}+y=x^2\log_ex,\brak{x>1}$. If $y\brak{2}=2$, then $y\brak{e}$ is equal to
    
    \hfill{[Jan 2023]}

		\begin{multicols}{4}
			\begin{enumerate}
				\item $\frac{4+e^2}{4}$
				\item $\frac{1+e^2}{4}$
				\item $\frac{2+e^2}{2}$
				\item $\frac{1+e^2}{2}$
			\end{enumerate}
		\end{multicols}

%18
    \item The number of $3$ digit numbers, that are divisible by either $3$ or $4$ but not divisible by $48$, is
    
    \hfill{[Jan 2023]}

        \begin{multicols}{4}
            \begin{enumerate}
                \item $472$
                \item $432$
                \item $507$
                \item $400$
            \end{enumerate}
        \end{multicols}

%19
    \item Let $R$ be a relation defined on $N$ as $a\ R\ b$ is $2a+3b$ is a multiple of 5, $a,b\in N$. Then $R$ is
    
    \hfill{[Jan 2023]}

		\begin{multicols}{2}
			\begin{enumerate}
				\item not reflexive
				\item transitive but not symmetric
				\item symmetric but not transitive
				\item an equivalence relation
			\end{enumerate}
		\end{multicols}

%20
    \item Consider a function $f:N\to R$, satisfying $f\brak{1}+2f\brak{2}+3f\brak{3}+\cdots+xf\brak{x}=x\brak{x+1}f\brak{x};x>2$ with $f\brak{1}=1$. Then $\frac{1}{f\brak{2022}}+\frac{1}{f\brak{2028}}$ is equal to 
    
    \hfill{[Jan 2023]}

		\begin{multicols}{4}
			\begin{enumerate}
				\item $8200$
				\item $8000$
				\item $8400$
				\item $8100$
			\end{enumerate}
		\end{multicols}
  

